\[ %%%%%%%%%%%%%%%%%%%%%%%%%%%%%%%%%%%%%%%%%%%%%%%%%%%%%%%%%%%%%%%%%%%%%%%%%%%%%%%% \subsection{Evaluation Metrics} \label{sec:evaluation_metrics} %%%%%%%%%%%%%%%%%%%%%%%%%%%%%%%%%%%%%%%%%%%%%%%%%%%%%%%%%%%%%%%%%%%%%%%%%%%%%%%% To comprehensively evaluate PHYSGEN, we employ a set of metrics measuring prediction accuracy, physical consistency, long-horizon stability, and computational efficiency. Unlike conventional forecasting benchmarks that focus only on prediction error, the evaluation of physics-guided models requires measuring whether generated trajectories remain physically meaningful over extended horizons. The complete evaluation framework is defined as: \begin{equation} \mathcal{M} = \{ \mathcal{M}_{accuracy}, \mathcal{M}_{physics}, \mathcal{M}_{stability}, \mathcal{M}_{efficiency} \}. \end{equation} %%%%%%%%%%%%%%%%%%%%%%%%%%%%%%%%%%%%%%%%%%%%%%%%%%%%%%%%%%%%%%%%%%%%%%%%%%%%%%%% \subsubsection{Prediction Accuracy Metrics} %%%%%%%%%%%%%%%%%%%%%%%%%%%%%%%%%%%%%%%%%%%%%%%%%%%%%%%%%%%%%%%%%%%%%%%%%%%%%%%% \paragraph{Mean Absolute Error (MAE).} The mean absolute error measures the average deviation between predicted and reference states: \begin{equation} MAE = \frac{1}{N} \sum_{i=1}^{N} | X_i-\hat{X}_i |. \end{equation} MAE provides an interpretable measure of average prediction accuracy. %%%%%%%%%%%%%%%%%%%%%%%%%%%%%%%%%%%%%%%%%%%%%%%%%%%%%%%%%%%%%%%%%%%%%%%%%%%%%%%% \paragraph{Root Mean Square Error (RMSE).} The root mean square error is defined as: \begin{equation} RMSE = \sqrt{ \frac{1}{N} \sum_{i=1}^{N} (X_i-\hat{X}_i)^2 }. \end{equation} RMSE emphasizes larger prediction deviations and is widely used for dynamical system forecasting. %%%%%%%%%%%%%%%%%%%%%%%%%%%%%%%%%%%%%%%%%%%%%%%%%%%%%%%%%%%%%%%%%%%%%%%%%%%%%%%% \paragraph{Relative $L_2$ Error.} For systems with different physical scales, the normalized error is computed: \begin{equation} E_{rel} = \frac{ \| X-\hat{X} \|_2 }{ \| X \|_2 }. \end{equation} This metric enables comparison between different physical systems. %%%%%%%%%%%%%%%%%%%%%%%%%%%%%%%%%%%%%%%%%%%%%%%%%%%%%%%%%%%%%%%%%%%%%%%%%%%%%%%% \subsubsection{Long-Horizon Forecasting Metrics} %%%%%%%%%%%%%%%%%%%%%%%%%%%%%%%%%%%%%%%%%%%%%%%%%%%%%%%%%%%%%%%%%%%%%%%%%%%%%%%% Since PHYSGEN focuses on long-term prediction, error evolution over time is explicitly evaluated. The horizon-dependent error is: \begin{equation} E(k) = \| X_{t+k} - \hat{X}_{t+k} \|. \end{equation} The average long-horizon error is: \begin{equation} LHE = \frac{1}{K} \sum_{k=1}^{K} E(k). \end{equation} This metric evaluates accumulated prediction degradation during recursive rollout. %%%%%%%%%%%%%%%%%%%%%%%%%%%%%%%%%%%%%%%%%%%%%%%%%%%%%%%%%%%%%%%%%%%%%%%%%%%%%%%% \paragraph{Stability Horizon.} The stability horizon measures the maximum prediction time before the error exceeds an acceptable threshold: \begin{equation} H_s = \max \{ k: E(k)<\epsilon \}. \end{equation} A larger $H_s$ indicates more stable long-term forecasting. %%%%%%%%%%%%%%%%%%%%%%%%%%%%%%%%%%%%%%%%%%%%%%%%%%%%%%%%%%%%%%%%%%%%%%%%%%%%%%%% \subsubsection{Physics Consistency Metrics} %%%%%%%%%%%%%%%%%%%%%%%%%%%%%%%%%%%%%%%%%%%%%%%%%%%%%%%%%%%%%%%%%%%%%%%%%%%%%%%% Prediction accuracy alone does not guarantee physical validity. Therefore, PHYSGEN evaluation includes physical constraint violations. %%%%%%%%%%%%%%%%%%%%%%%%%%%%%%%%%%%%%%%%%%%%%%%%%%%%%%%%%%%%%%%%%%%%%%%%%%%%%%%% \paragraph{Physics Residual Error.} For systems governed by: \begin{equation} \mathcal{F}(X)=0, \end{equation} the physics residual is: \begin{equation} E_{phys} = \| \mathcal{F}(\hat{X}) \|^2. \end{equation} Lower values indicate better physical consistency. %%%%%%%%%%%%%%%%%%%%%%%%%%%%%%%%%%%%%%%%%%%%%%%%%%%%%%%%%%%%%%%%%%%%%%%%%%%%%%%% \paragraph{Conservation Error.} For conserved quantities $I(X)$: \begin{equation} E_{cons} = \| I(\hat{X}_{t+k}) - I(X_t) \|. \end{equation} Examples include: \begin{itemize} \item mass conservation; \item energy conservation; \item momentum preservation. \end{itemize} %%%%%%%%%%%%%%%%%%%%%%%%%%%%%%%%%%%%%%%%%%%%%%%%%%%%%%%%%%%%%%%%%%%%%%%%%%%%%%%% \subsubsection{Trajectory Stability Metrics} %%%%%%%%%%%%%%%%%%%%%%%%%%%%%%%%%%%%%%%%%%%%%%%%%%%%%%%%%%%%%%%%%%%%%%%%%%%%%%%% To evaluate dynamical stability, we measure trajectory divergence: \begin{equation} D_k = \| \hat{X}_{t+k} - X_{t+k} \|. \end{equation} The divergence growth rate is estimated as: \begin{equation} \lambda_{eff} = \frac{1}{K} \log \frac{D_K}{D_0}. \end{equation} A smaller effective growth rate indicates improved stability. %%%%%%%%%%%%%%%%%%%%%%%%%%%%%%%%%%%%%%%%%%%%%%%%%%%%%%%%%%%%%%%%%%%%%%%%%%%%%%%% \paragraph{Latent Stability Score.} The evolution of latent states is measured by: \begin{equation} S_{latent} = \frac{1}{K} \sum_{k=1}^{K} \| H_{t+k}-H_{t+k-1} \|. \end{equation} This evaluates whether the learned latent dynamics remain bounded. %%%%%%%%%%%%%%%%%%%%%%%%%%%%%%%%%%%%%%%%%%%%%%%%%%%%%%%%%%%%%%%%%%%%%%%%%%%%%%%% \subsubsection{Graph Structure Metrics} %%%%%%%%%%%%%%%%%%%%%%%%%%%%%%%%%%%%%%%%%%%%%%%%%%%%%%%%%%%%%%%%%%%%%%%%%%%%%%%% Since PHYSGEN models physical systems as evolving graphs, graph consistency is also evaluated. The topology error is defined as: \begin{equation} E_{graph} = D(G_t,\hat{G}_t), \end{equation} where $D$ measures differences in: \begin{itemize} \item connectivity; \item edge weights; \item spatial relationships. \end{itemize} %%%%%%%%%%%%%%%%%%%%%%%%%%%%%%%%%%%%%%%%%%%%%%%%%%%%%%%%%%%%%%%%%%%%%%%%%%%%%%%% \subsubsection{Computational Metrics} %%%%%%%%%%%%%%%%%%%%%%%%%%%%%%%%%%%%%%%%%%%%%%%%%%%%%%%%%%%%%%%%%%%%%%%%%%%%%%%% In addition to predictive performance, computational efficiency is evaluated. The following metrics are measured: \begin{itemize} \item training time; \item inference time per prediction step; \item GPU memory consumption; \item number of trainable parameters. \end{itemize} The computational complexity is reported as: \begin{equation} C_{model} = O(|E|d^2), \end{equation} where $|E|$ is the number of graph edges and $d$ is the hidden dimension. %%%%%%%%%%%%%%%%%%%%%%%%%%%%%%%%%%%%%%%%%%%%%%%%%%%%%%%%%%%%%%%%%%%%%%%%%%%%%%%% \subsubsection{Composite PHYSGEN Score} %%%%%%%%%%%%%%%%%%%%%%%%%%%%%%%%%%%%%%%%%%%%%%%%%%%%%%%%%%%%%%%%%%%%%%%%%%%%%%%% To provide a unified evaluation, a composite score can be defined: \begin{equation} S_{PHYSGEN} = w_aS_{acc} + w_pS_{phys} + w_sS_{stab}, \end{equation} where: \begin{itemize} \item $S_{acc}$ measures prediction accuracy; \item $S_{phys}$ measures physical consistency; \item $S_{stab}$ measures long-horizon stability. \end{itemize} The weights $w_a,w_p,w_s$ are selected according to the target application. %%%%%%%%%%%%%%%%%%%%%%%%%%%%%%%%%%%%%%%%%%%%%%%%%%%%%%%%%%%%%%%%%%%%%%%%%%%%%%%% \subsubsection{Summary} %%%%%%%%%%%%%%%%%%%%%%%%%%%%%%%%%%%%%%%%%%%%%%%%%%%%%%%%%%%%%%%%%%%%%%%%%%%%%%%% The evaluation framework for PHYSGEN extends beyond conventional prediction metrics by incorporating physical validity and dynamical stability. This multi-dimensional evaluation protocol enables rigorous assessment of physics-guided learning methods for long-horizon prediction of nonlinear physical systems. \] 
